\chapter{Evaluarea performanțelor}
\label{ch:evaluare}

\section{Metodologie}

\paragraph{Mediu.} Toate componentele rulează pe aceeași mașină, într-un VM Docker Desktop cu
\textbf{5 vCPU} (Apple Silicon, arm64). Generatorul de încărcare, nginx, cele două instanțe, nodurile de
cache și bazele de date concurează pentru aceleași nuclee, deci rezultatele absolute sunt specifice
acestui mediu. Comparațiile între configurații sunt însă relevante, pentru că au fost obținute în
condiții identice.

\paragraph{Sarcina de lucru.} Generatorul \texttt{loadgen} rulează în rețeaua Docker și trimite cereri
prin nginx. Implicit folosește 64 de clienți \emph{closed-loop} (fiecare trimite cererea următoare
imediat după răspuns), 95\% citiri (din care 10\% clasamente de categorie), 5\% scrieri de preț și
popularitate Zipf cu $s = 0{,}99$ peste 100.000 de produse și 200 de categorii. Fiecare rulare are 5~s
de încălzire (nemăsurate) și 15~s de măsurare. Contoarele aplicației sunt resetate după încălzire.
Schimbarea modului golește cache-ul.

\paragraph{Metrici.} Debitul (cereri reușite/s), latențele p50/p95/p99 (histograme HDR, rezoluție 1~µs,
3 cifre semnificative), rata de hit (din contoarele aplicației, însumate pe instanțe), numărul de
interogări și tranzacții în baza de date. În E8, latența este măsurată de la momentul \emph{programat}
al cererii, ca să evităm \emph{coordinated omission} \cite{tene}.

\begin{table}[H]
\centering\small
\begin{tabularx}{\linewidth}{llX}
\toprule
Exp. & Variabilă & Întrebare \\
\midrule
E1 & backend (6) & Ce câștig aduce fiecare tip de cache față de baza de date? \\
E2 & exponent Zipf & Cum depinde rata de hit de asimetria accesărilor, cu memorie limitată? \\
E3 & proporția de scrieri × strategie & Ce costă coerența la scriere? \\
E4 & clienți concurenți (1–256) & Cum scalează sistemul cu concurența? (1 client = secvențial) \\
E5 & instanțe web (1–2) & Câștig din scalarea orizontală? \\
E6 & noduri de cache (1–3) & Crește rata de hit cu memoria agregată a clusterului? \\
E7 & single-flight on/off & Cât reduce protecția la stampede încărcarea bazei? \\
E8 & backend, debit fix & Cât CPU costă o cerere în fiecare arhitectură? \\
SIM & politică × capacitate & Cât de aproape de optim sunt LRU/FIFO/CLOCK? \\
ENG & fire × shard-uri & Cât de bine se paralelizează motorul \texttt{memtable}? \\
\bottomrule
\end{tabularx}
\caption{Experimentele realizate.}
\end{table}

\section{E1 — Comparația backend-urilor}

\begin{figure}[H]
  \centering% E1 — comparația backend-urilor.
\begin{tikzpicture}
\begin{groupplot}[group style={group size=2 by 1, horizontal sep=17mm},
    width=0.5\linewidth, height=55mm, ybar, xtick=data,
    symbolic x coords={none,local,cluster,tiered,redis,pgmem},
    x tick label style={rotate=35, anchor=east, font=\footnotesize}, enlarge x limits=0.1]
  \nextgroupplot[ylabel={cereri / s}, bar width=8pt, ymin=0,
      nodes near coords, nodes near coords style={font=\tiny, rotate=90, anchor=west},
      point meta={\thisrow{rps}/1000}, nodes near coords={\pgfmathprintnumber[fixed,precision=1]{\pgfplotspointmeta}k},
      kticks, ymax=88000,
      title={\small (a) debit}]
    \addplot[fill=c1!70, draw=c1] table[x=backend, y=rps, col sep=comma] {\datadir e1.csv};
  \nextgroupplot[ylabel={latență (ms)}, ymode=log, bar width=5pt, log origin=infty,
      legend style={at={(0.02,0.98)}, anchor=north west}, title={\small (b) latență}]
    \addplot[fill=c3!70, draw=c3] table[x=backend, y=p50, col sep=comma] {\datadir e1.csv};
    \addplot[fill=c2!70, draw=c2] table[x=backend, y=p95, col sep=comma] {\datadir e1.csv};
    \addplot[fill=c5!70, draw=c5] table[x=backend, y=p99, col sep=comma] {\datadir e1.csv};
    \legend{p50, p95, p99}
\end{groupplot}
\end{tikzpicture}

  \caption{Debitul și latența pentru cele șase backend-uri (64 de clienți, 95\% citiri, Zipf 0,99).}
  \label{fig:e1}
\end{figure}

\begin{table}[H]
\centering\small
\begin{tabular}{lrrrr}
\toprule
Backend & cereri/s & p50 (ms) & p99 (ms) & hit (\%) \\
\midrule
\texttt{none}    & \R{e1/none/rps}    & \R{e1/none/p50}    & \R{e1/none/p99}    & — \\
\texttt{local}   & \R{e1/local/rps}   & \R{e1/local/p50}   & \R{e1/local/p99}   & \R{e1/local/hit} \\
\texttt{cluster} & \R{e1/cluster/rps} & \R{e1/cluster/p50} & \R{e1/cluster/p99} & \R{e1/cluster/hit} \\
\texttt{tiered}  & \R{e1/tiered/rps}  & \R{e1/tiered/p50}  & \R{e1/tiered/p99}  & \R{e1/tiered/hit} \\
\texttt{redis}   & \R{e1/redis/rps}   & \R{e1/redis/p50}   & \R{e1/redis/p99}   & \R{e1/redis/hit} \\
\texttt{pgmem}   & \R{e1/pgmem/rps}   & \R{e1/pgmem/p50}   & \R{e1/pgmem/p99}   & \R{e1/pgmem/hit} \\
\bottomrule
\end{tabular}
\caption{Rezultatele E1.}
\label{tab:e1}
\end{table}

\paragraph{Interpretare.} Toate cache-urile in-memory reale depășesc baza de date. Cache-ul local
are cel mai mare debit (\R{e1/local/rps} cereri/s, de \R{e1/local/speedup}× față de \texttt{none}) și cea mai
mică latență mediană (\R{e1/local/p50}~ms), pentru că un hit nu părăsește procesul. Rata lui de hit
(\R{e1/local/hit}\%) este însă mai mică decât a cache-urilor partajate (\R{e1/cluster/hit}\% pentru cluster,
\R{e1/redis/hit}\% pentru Redis). Cauza: fiecare dintre cele două instanțe își încălzește separat propria copie,
iar notificările \texttt{NOTIFY} invalidează ambele copii la fiecare scriere. Clusterul propriu
(\R{e1/cluster/rps} cereri/s) este comparabil cu Redis (\R{e1/redis/rps}), deși folosește un protocol și un
motor scrise de la zero.

Varianta \texttt{tiered} (\R{e1/tiered/rps}) are rata de hit cea mai mare, dar debit mai mic decât clusterul
simplu. Având un nivel local, ea reactivează \texttt{NOTIFY}, care serializează commit-urile
(\cref{sec:optimizari}), și dublează scrierile în cache la fiecare miss. Câștigul nivelului L1 apare doar
când latența rețelei până la L2 e mare, ceea ce nu e cazul într-o rețea Docker locală.

Rezultatul cel mai instructiv este \texttt{pgmem}: tabela in-memory din PostgreSQL este \emph{mai lentă} decât
baza de date fără cache (\R{e1/pgmem/rps} față de \R{e1/none/rps} cereri/s), deși are rată de hit
\R{e1/pgmem/hit}\%. Un hit costă aici la fel de mult ca interogarea originală: aceeași parsare, aceeași
execuție SQL și același pool de conexiuni. După optimizările din \cref{sec:optimizari}, interogarea sursă
este un simplu lookup pe cheia primară, servit din \emph{shared buffers}. Un miss costă chiar trei operații
(citire cache, citire DB, \emph{upsert}). Concluzia: \textbf{o tabelă in-memory într-un SGBD nu este un cache
eficient pentru date pe care SGBD-ul le servește deja din memorie}. Valoarea ei apare doar când rezultatul
cache-uit este mult mai scump decât un lookup.


\section{E2 — Asimetria accesărilor}

\begin{figure}[H]
  \centering% E2 — asimetria Zipf (cache = 10% din chei, doar citiri).
\begin{tikzpicture}
\begin{groupplot}[group style={group size=2 by 1, horizontal sep=17mm},
    width=0.5\linewidth, height=55mm, xlabel={exponent Zipf $s$}, xtick={0.6,0.8,0.99,1.2}]
  \nextgroupplot[ylabel={rata de hit}, ymin=0, ymax=1, legend pos=south east, title={\small (a) rata de hit}]
    \addplot table[x=label, y=local_hit_ratio, col sep=comma] {\datadir e2.csv};
    \addplot table[x=label, y=cluster_hit_ratio, col sep=comma] {\datadir e2.csv};
    \addplot[c6, thick, dashed, mark=o] table[x=label, y=ideal, col sep=comma] {\datadir e2.csv};
    \legend{local, cluster, ideal IRM}
  \nextgroupplot[ylabel={cereri / s}, ymin=0, legend pos=north west, title={\small (b) debit},
      kticks]
    \addplot[c6, mark=o, thick] table[x=label, y=none_throughput_rps, col sep=comma] {\datadir e2.csv};
    \addplot table[x=label, y=local_throughput_rps, col sep=comma] {\datadir e2.csv};
    \addplot table[x=label, y=cluster_throughput_rps, col sep=comma] {\datadir e2.csv};
    \legend{none, local, cluster}
\end{groupplot}
\end{tikzpicture}

  \caption{Rata de hit și debitul în funcție de exponentul Zipf. Cache-ul are 10\% din chei, iar
           traficul conține doar citiri. Linia punctată este limita teoretică IRM (\cref{eq:ideal}).}
  \label{fig:e2}
\end{figure}

\paragraph{Interpretare.} Cu memoria limitată la 10\% din chei, rata de hit crește puternic cu asimetria:
de la \R{e2/local/0.6/hit_ratio}\% la $s=0{,}6$ la \R{e2/local/1.2/hit_ratio}\% la $s=1{,}2$. Valorile urmăresc
îndeaproape limita teoretică IRM (\cref{eq:ideal}): diferența de 2–8 puncte reprezintă costul unei politici
\emph{online} (LRU) față de cache-ul static optim. Cele două backend-uri au practic aceeași rată de hit, pentru
că au aceeași capacitate totală.

Debitul cache-ului local se dublează pe intervalul studiat (\R{e2/local/0.6/throughput_rps} $\rightarrow$
\R{e2/local/1.2/throughput_rps} cereri/s), în timp ce baza de date rămâne aproape constantă. La asimetrie mică
($s=0{,}6$), clusterul este însă \emph{mai lent} decât baza de date (\R{e2/cluster/0.6/throughput_rps} față de
\R{e2/none/0.6/throughput_rps}). Explicația vine din \cref{eq:amat}: un cache la distanță adaugă o ședere în
rețea fiecărei cereri, deci câștigă doar dacă $h > T_{\text{cache}}/T_{\text{DB}}$. Când interogarea sursă este
ieftină, pragul de rentabilitate este ridicat, iar cu doar \R{e2/cluster/0.6/hit_ratio}\% hit-uri el nu este atins.


\section{E3 — Strategii de scriere}

\begin{figure}[H]
  \centering% E3 — strategii de scriere.
\begin{tikzpicture}
\begin{groupplot}[group style={group size=2 by 1, horizontal sep=17mm},
    width=0.5\linewidth, height=55mm, xlabel={proporția de citiri}, xtick={0.5,0.8,0.95}]
  \nextgroupplot[ylabel={cereri / s}, ymin=0, legend pos=south east, title={\small (a) debit},
      kticks]
    \addplot table[x=label, y=invalidate_throughput_rps, col sep=comma] {\datadir e3.csv};
    \addplot table[x=label, y=write-through_throughput_rps, col sep=comma] {\datadir e3.csv};
    \addplot table[x=label, y=write-behind_throughput_rps, col sep=comma] {\datadir e3.csv};
    \legend{invalidate, write-through, write-behind}
  \nextgroupplot[ylabel={tranzacții DB / 1000 scrieri}, ymode=log, ymin=150, ymax=1300, ytick={200,300,500,1000}, log ticks with fixed point,
      title={\small (b) scrieri în DB}]
    \addplot table[x=label, y=invalidate_tx_per_1k, col sep=comma] {\datadir e3.csv};
    \addplot table[x=label, y=write-through_tx_per_1k, col sep=comma] {\datadir e3.csv};
    \addplot table[x=label, y=write-behind_tx_per_1k, col sep=comma] {\datadir e3.csv};
\end{groupplot}
\end{tikzpicture}

  \caption{Debitul și numărul de tranzacții în baza de date la 1.000 de scrieri, pentru cele trei
           strategii (backend \texttt{cluster}).}
  \label{fig:e3}
\end{figure}

\paragraph{Interpretare.} Diferențele apar la proporții mari de scrieri. La 50\% scrieri, invalidarea păstrează
doar \R{e3/invalidate/0.5/hit_ratio}\% hit-uri (fiecare scriere șterge o cheie fierbinte), iar debitul este
\R{e3/invalidate/0.5/throughput_rps} cereri/s. Write-through ridică rata de hit la
\R{e3/write-through/0.5/hit_ratio}\%. Write-behind atinge \R{e3/write-behind/0.5/throughput_rps} cereri/s
(\R{e3/write-behind/0.5/hit_ratio}\% hit-uri) și reduce numărul de tranzacții la \R{e3/write-behind/0.5/tx1k} per
1.000 de scrieri, față de $\approx$1.000 la strategiile sincrone. Câștigul vine din grupare (mai multe
modificări într-un \texttt{UPDATE}) și din combinarea modificărilor repetate ale aceluiași produs fierbinte.

La 95\% citiri strategiile converg, iar write-behind este chiar puțin mai lent. Loturile sunt mici (rata de
scriere e redusă), deci gruparea aduce puțin, iar fiecare scriere plătește un \emph{read-modify-write} în cache.

\paragraph{Lecție de proiectare: backlog-ul.} O primă versiune folosea un singur worker per instanță. Debitul
măsurat era mai mare, dar la finalul rulării erau $\approx$181.000 de scrieri încă în coadă, golite în
$\approx$5~s după oprirea încărcării. Sistemul confirma deci scrieri mai repede decât le putea persista. Am
împărțit write-behind în 4 benzi paralele, partiționate după \texttt{id mod 4}, ceea ce păstrează ordinea
modificărilor fiecărui produs. Cu ele coada rămâne goală la final, iar cifrele de mai sus reprezintă un debit
\emph{sustenabil}. Fereastra de inconsistență dintre cache și baza de date este de ordinul unui lot (milisecunde).


\section{E4 — Scalabilitatea cu numărul de clienți}

\begin{figure}[H]
  \centering% E4 — scalabilitatea cu numărul de clienți.
\begin{tikzpicture}
\begin{groupplot}[group style={group size=2 by 1, horizontal sep=17mm},
    width=0.5\linewidth, height=55mm, xlabel={clienți concurenți}, xmode=log, log basis x=2,
    xtick={1,4,16,64,256}, xticklabels={1,4,16,64,256}]
  \nextgroupplot[ylabel={cereri / s}, ymin=0, legend pos=north west, title={\small (a) debit},
      kticks]
    \addplot[c6, mark=o, thick] table[x=label, y=none_throughput_rps, col sep=comma] {\datadir e4.csv};
    \addplot table[x=label, y=local_throughput_rps, col sep=comma] {\datadir e4.csv};
    \addplot table[x=label, y=cluster_throughput_rps, col sep=comma] {\datadir e4.csv};
    \legend{none, local, cluster}
  \nextgroupplot[ylabel={latență p99 (ms)}, ymode=log, legend pos=north west, title={\small (b) latență p99}]
    \addplot[c6, mark=o, thick] table[x=label, y=none_p99_ms, col sep=comma] {\datadir e4.csv};
    \addplot table[x=label, y=local_p99_ms, col sep=comma] {\datadir e4.csv};
    \addplot table[x=label, y=cluster_p99_ms, col sep=comma] {\datadir e4.csv};
\end{groupplot}
\end{tikzpicture}

  \caption{Debitul și latența p99 în funcție de numărul de clienți concurenți. Un client corespunde
           execuției secvențiale.}
  \label{fig:e4}
\end{figure}

\paragraph{Interpretare.} Cu un singur client, adică execuția secvențială, sistemul servește
\R{e4/none/1/throughput_rps} cereri/s fără cache și \R{e4/local/1/throughput_rps} cu cache local, la o latență
mediană de numai \R{e4/local/1/p50_ms}~ms. Debitul crește aproape liniar cu concurența, conform legii lui Little
($\lambda = L/W$ cu $W$ aproape constant), până la saturarea CPU-ului. Maximul este de
\R{e4/local/64/throughput_rps} cereri/s pentru cache-ul local la 64 de clienți (de $\approx$9× față de
secvențial) și de \R{e4/none/128/throughput_rps} pentru baza de date la 128 de clienți.

După saturare, debitul nu mai crește, iar clienții suplimentari doar așteaptă la coadă: latența p99 crește
proporțional cu concurența (\R{e4/none/256/p99_ms}~ms fără cache la 256 de clienți). Debitul chiar scade ușor
din cauza contenției (comutări de context, lock-uri în PostgreSQL). De la 16 clienți în sus, ambele cache-uri
păstrează p99 sub varianta fără cache. La concurență mică, clusterul are p99 ușor mai mare
(\R{e4/cluster/1/p99_ms} față de \R{e4/none/1/p99_ms}~ms la un client), pentru că adaugă o ședere în rețea.


\section{E5–E7 — Scale-out, memorie agregată, stampede}

\begin{figure}[H]
  \centering\resizebox{\linewidth}{!}{% E5 scale-out, E6 noduri, E7 stampede.
\begin{tikzpicture}
\begin{groupplot}[group style={group size=3 by 1, horizontal sep=21mm}, width=0.36\linewidth, height=52mm]
  \nextgroupplot[ybar, bar width=9pt, ymin=0, xtick={1,2}, xlabel={instanțe web}, ylabel={cereri / s},
      enlarge x limits=0.6, ymax=108000, legend columns=2, legend style={at={(0.5,0.98)}, anchor=north},
      kticks, title={\small (a) E5: scale-out}]
    \addplot[fill=c6!50, draw=c6] table[x=label, y=none_throughput_rps, col sep=comma] {\datadir e5.csv};
    \addplot[fill=c2!70, draw=c2] table[x=label, y=cluster_throughput_rps, col sep=comma] {\datadir e5.csv};
    \legend{none, cluster}
  \nextgroupplot[xtick={1,2,3}, xlabel={noduri de cache}, ylabel={rata de hit}, ymin=0, ymax=1,
      title={\small (b) E6: memorie agregată}]
    \addplot table[x=label, y=cluster_hit_ratio, col sep=comma] {\datadir e6.csv};
  \nextgroupplot[ybar, bar width=16pt, ymin=0, xtick={0,1}, xticklabels={fără, cu}, xlabel={single-flight},
      ylabel={interogări DB}, enlarge x limits=0.6, scaled y ticks=false, ymax=170000,
      yticklabel={\pgfmathprintnumber[fixed]{\tick}}, title={\small (c) E7: stampede}]
    \addplot[fill=c5!60, draw=c5] table[x=coalesce, y=db_reads, col sep=comma] {\datadir e7.csv};
\end{groupplot}
\end{tikzpicture}
}
  \caption{(a) Una vs. două instanțe web (128 de clienți); (b) rata de hit cu 1–3 noduri de câte 10.000
           de intrări (doar citiri); (c) interogări DB la pornirea la rece, cu Zipf 1,4 și 128 de clienți.}
  \label{fig:e567}
\end{figure}

\paragraph{E5 — scale-out.} Pe o singură mașină cu 5 vCPU, a doua instanță web \emph{nu} crește debitul
(\R{e5/cluster/1/throughput_rps} cereri/s cu o instanță, \R{e5/cluster/2/throughput_rps} cu două, backend
cluster). O singură instanță tokio folosește deja toate nucleele, iar a doua concurează pentru aceleași
resurse și adaugă conexiuni și coordonare. Rezultatul ilustrează legea lui Amdahl la nivel de sistem: scalarea
orizontală aduce debit doar dacă adaugă resurse noi (mașini). Pe același hardware, beneficiul celei de-a doua
instanțe este disponibilitatea: nginx redirecționează traficul dacă o instanță cade.

\paragraph{E6 — memorie agregată.} Cu 10.000 de intrări pe nod, rata de hit crește de la
\R{e6/cluster/1/hit_ratio}\% (1 nod) la \R{e6/cluster/2/hit_ratio}\% (2 noduri) și \R{e6/cluster/3/hit_ratio}\%
(3 noduri), aproape de limitele teoretice pentru 10/20/30\% din chei
(\R{sim/zipf/0.99/10/ideal}, \R{sim/zipf/0.99/20/ideal}, \R{sim/zipf/0.99/30/ideal}\%). Partiționarea prin
consistent hashing face ca memoria nodurilor să se adune, nu să se dubleze. Debitul rămâne constant
($\approx$\R{e6/cluster/3/throughput_rps} cereri/s), pentru că la doar citiri gâtul de sticlă este CPU-ul comun.

\paragraph{E7 — cache stampede.} La pornirea la rece, cu o distribuție foarte asimetrică ($s = 1{,}4$),
single-flight reduce interogările către baza de date de la \R{e7/false/db_reads} la \R{e7/true/db_reads}
($-$29\%): \R{e7/true/coalesced} de cereri au fost servite de o interogare aflată deja în curs. Debitul crește de la
\R{e7/false/throughput_rps} la \R{e7/true/throughput_rps} cereri/s.


\section{E8 — Costul CPU la debit fix}

\begin{figure}[H]
  \centering% E8 — CPU consumat la debit fix (8.000 cereri/s), pe roluri.
\begin{tikzpicture}
\begin{axis}[width=0.8\linewidth, height=60mm, ybar stacked, bar width=14pt, ymin=0,
    symbolic x coords={none,local,cluster,tiered,redis,pgmem}, xtick=data,
    ylabel={CPU (\% dintr-un nucleu)}, legend style={at={(1.02,1)}, anchor=north west},
    enlarge x limits=0.12, reverse legend]
  \addplot[fill=c3!70, draw=c3] table[x=backend, y=app, col sep=comma] {\datadir e8.csv};
  \addplot[fill=c1!70, draw=c1] table[x=backend, y=db, col sep=comma] {\datadir e8.csv};
  \addplot[fill=c2!70, draw=c2] table[x=backend, y=cache, col sep=comma] {\datadir e8.csv};
  \addplot[fill=c6!35, draw=c6] table[x=backend, y=nginx, col sep=comma] {\datadir e8.csv};
  \legend{aplicație, PostgreSQL, servere de cache, nginx}
\end{axis}
\end{tikzpicture}

  \caption{CPU consumat de fiecare rol la un debit fix de 8.000 de cereri/s (open-loop). 100\% = un
           nucleu complet.}
  \label{fig:e8}
\end{figure}

\paragraph{Interpretare.} La același debit (8.000 de cereri/s, toate servite fără erori), suma CPU consumată
de aplicație, baza de date și serverele de cache este de \R{e8/none/usreq}~µs per cerere fără cache și
\R{e8/local/usreq}~µs cu cache local ($-$27\%). Pentru cluster este \R{e8/cluster/usreq}~µs, iar pentru Redis
\R{e8/redis/usreq}~µs. Consumul PostgreSQL scade de la \R{e8/none/db}\% la \R{e8/cluster/db}\% dintr-un nucleu
cu cache-ul partajat. Cu cache local scade doar la \R{e8/local/db}\%, din cauza \texttt{NOTIFY} și a ratei de hit
mai mici. Varianta \texttt{pgmem} este cea mai scumpă (\R{e8/pgmem/usreq}~µs per cerere): CPU-ul economisit în
baza principală se mută, amplificat, în serverul PostgreSQL care ține cache-ul. Măsurătoarea confirmă
concluzia din E1 dintr-un unghi independent de saturare.


\section{Politici de evacuare (simulare)}
\label{sec:sim}

\begin{figure}[H]
  \centering% Simulare politici de evacuare.
\begin{tikzpicture}
\begin{groupplot}[group style={group size=2 by 2, horizontal sep=17mm, vertical sep=18mm},
    width=0.5\linewidth, height=50mm, xmode=log, xtick={1,2,5,10,20,50}, xticklabels={1,2,5,10,20,50},
    xlabel={capacitate (\% din chei)}, ylabel={rata de hit}, ymin=0, ymax=1]
  \nextgroupplot[title={\small Zipf $s=0{,}8$}, legend pos=north west]
    \addplot table[x=cap, y=lru, col sep=comma] {\datadir sim_zipf_08.csv};
    \addplot table[x=cap, y=fifo, col sep=comma] {\datadir sim_zipf_08.csv};
    \addplot table[x=cap, y=clock, col sep=comma] {\datadir sim_zipf_08.csv};
    \addplot[c6, thick, dashed] table[x=cap, y=ideal, col sep=comma] {\datadir sim_zipf_08.csv};
    \legend{LRU, FIFO, CLOCK, ideal}
  \nextgroupplot[title={\small Zipf $s=0{,}99$}]
    \addplot table[x=cap, y=lru, col sep=comma] {\datadir sim_zipf_099.csv};
    \addplot table[x=cap, y=fifo, col sep=comma] {\datadir sim_zipf_099.csv};
    \addplot table[x=cap, y=clock, col sep=comma] {\datadir sim_zipf_099.csv};
    \addplot[c6, thick, dashed] table[x=cap, y=ideal, col sep=comma] {\datadir sim_zipf_099.csv};
  \nextgroupplot[title={\small Zipf $s=0{,}8$ + 20\% scanări}]
    \addplot table[x=cap, y=lru, col sep=comma] {\datadir sim_zipf_scan_08.csv};
    \addplot table[x=cap, y=fifo, col sep=comma] {\datadir sim_zipf_scan_08.csv};
    \addplot table[x=cap, y=clock, col sep=comma] {\datadir sim_zipf_scan_08.csv};
  \nextgroupplot[title={\small Zipf $s=0{,}99$ + 20\% scanări}]
    \addplot table[x=cap, y=lru, col sep=comma] {\datadir sim_zipf_scan_099.csv};
    \addplot table[x=cap, y=fifo, col sep=comma] {\datadir sim_zipf_scan_099.csv};
    \addplot table[x=cap, y=clock, col sep=comma] {\datadir sim_zipf_scan_099.csv};
\end{groupplot}
\end{tikzpicture}

  \caption{Rata de hit în funcție de capacitate (2 milioane de cereri, 100.000 de chei; prima jumătate
           doar încălzește cache-ul). Rândul de jos adaugă 20\% scanări secvențiale peste chei reci.}
  \label{fig:sim}
\end{figure}

\paragraph{Interpretare.} Pentru toate capacitățile, ordinea este aceeași: CLOCK $\gtrsim$ LRU $>$ FIFO. La
$s = 0{,}99$ și 10\% capacitate: CLOCK \R{sim/zipf/0.99/10/clock}\%, LRU \R{sim/zipf/0.99/10/lru}\%, FIFO
\R{sim/zipf/0.99/10/fifo}\%, față de limita teoretică \R{sim/zipf/0.99/10/ideal}\%. Faptul că CLOCK egalează sau
depășește LRU este un rezultat cunoscut. Sub modelul IRM recența nu este semnalul optim, iar „a doua șansă"
acordată intrărilor referențiate acționează ca un filtru slab de frecvență. CLOCK este în plus mai ieftin,
pentru că un hit nu modifică lista.

Scanările secvențiale (rândul de jos) degradează toate politicile: la $s=0{,}99$ și 10\% capacitate, LRU
scade de la \R{sim/zipf/0.99/10/lru}\% la \R{sim/zipf+scan/0.99/10/lru}\%. Cheile reci, accesate o singură dată,
împing din cache setul fierbinte. Diferența de 5–14 puncte față de optim și sensibilitatea la scanări motivează
politicile cu admitere bazată pe frecvență (W-TinyLFU, ARC).


\section{Scalabilitatea motorului \texttt{memtable}}

\begin{figure}[H]
  \centering\resizebox{\linewidth}{!}{% Scalabilitatea motorului memtable.
\begin{tikzpicture}
\begin{groupplot}[group style={group size=2 by 1, horizontal sep=17mm},
    width=0.5\linewidth, height=58mm, xlabel={fire}, xtick={1,2,4,5,8}]
  \nextgroupplot[ylabel={Mops / s}, ymin=0, ymax=19, title={\small (a) debit (LRU, Zipf 0{,}99)},
      legend style={font=\scriptsize, at={(0.02,0.98)}, anchor=north west}, legend columns=2]
    \addplot table[x=threads, y=s1_mops, col sep=comma] {\datadir engine.csv};
    \addplot table[x=threads, y=s4_mops, col sep=comma] {\datadir engine.csv};
    \addplot table[x=threads, y=s16_mops, col sep=comma] {\datadir engine.csv};
    \addplot table[x=threads, y=s64_mops, col sep=comma] {\datadir engine.csv};
    \legend{1 shard, 4 shard-uri, 16 shard-uri, 64 shard-uri}
  \nextgroupplot[ylabel={accelerare $S(p)$}, ymin=0, ymax=8, title={\small (b) accelerare, 64 shard-uri},
      legend style={font=\scriptsize, at={(0.02,0.98)}, anchor=north west}]
    \addplot table[x=threads, y=s64_speedup, col sep=comma] {\datadir engine.csv};
    \addplot table[x=threads, y=s64clock_speedup, col sep=comma] {\datadir engine.csv};
    \addplot table[x=threads, y=s64unif_speedup, col sep=comma] {\datadir engine.csv};
    \addplot[c5, mark=pentagon*, thick] table[x=threads, y=s1_speedup, col sep=comma] {\datadir engine.csv};
    \addplot[c6, dashed, thick, domain=1:8, samples=9, mark=none] {min(x,5)};
    \legend{{LRU, Zipf 0{,}99}, {CLOCK, Zipf 0{,}99}, LRU uniform, 1 shard (LRU), ideal}
\end{groupplot}
\end{tikzpicture}
}
  \caption{Debitul motorului (90\% citiri, Zipf 0,99, valori de 256 B) și accelerarea față de un fir, în
           funcție de numărul de fire și de shard-uri. Linia punctată este accelerarea ideală, limitată de
           cele 5 nuclee.}
  \label{fig:engine}
\end{figure}

\paragraph{Interpretare.} Pe un singur fir, motorul execută \R{eng/s64/1/mops} milioane de operații pe
secundă ($\approx$100~ns per operație, incluzând dispersia cheii și blocarea mutex-ului). Cu \textbf{un singur
shard}, adăugarea de fire \emph{reduce} debitul la \R{eng/s1/5/mops}~Mops la 5 fire. Toate firele se
serializează pe același mutex, iar linia de cache a lock-ului „sare" între nuclee: este cazul extrem din
legea lui Amdahl, cu partea serială $\approx 100\%$. Sharding-ul elimină acest efect: cu 64 de shard-uri,
debitul crește la \R{eng/s64/5/mops}~Mops.

Accelerarea rămâne însă modestă cu acces Zipf (\R{eng/s64/5/speedup}× la 5 fire), din cauza
\textbf{cheilor fierbinți}. La $s = 0{,}99$, cea mai populară cheie primește $\approx$8\% din toate accesările,
deci shard-ul ei și contorul de referințe al valorii ei devin puncte de serializare. Cu acces uniform,
accelerarea urcă la \R{eng/s64unif/5/speedup}×. Echivalentul Amdahl este o fracțiune paralelizabilă de
$\approx$75\% pentru acces uniform, față de $\approx$38\% pentru Zipf. Paradoxal, un fir singur este de două ori
mai rapid cu Zipf decât cu acces uniform (\R{eng/s64/1/mops} față de \R{eng/s64unif/1/mops}~Mops): setul
fierbinte încape în cache-urile procesorului. CLOCK câștigă până la 14\% față de LRU pe un fir (\R{eng/s64clock/1/mops} față de
\R{eng/s64/1/mops}~Mops), pentru că o citire nu mai mută noduri în listă. Peste 5 fire nu există câștig suplimentar, fiindcă VM-ul are 5 nuclee.
Problema cheilor fierbinți este cunoscută în sistemele reale \cite{memcache-fb} și se rezolvă prin replicarea
lor sau printr-un L1 per fir.


\section{Amenințări la validitate}

\begin{itemize}
  \item \textbf{Colocare}: generatorul și sistemul testat împart aceleași 5 nuclee, deci sub saturație o
        parte din CPU este consumată de generator și de nginx. Debitele absolute sunt astfel subestimate
        față de o implementare cu mașini separate.
  \item \textbf{Rețea virtuală}: latențele de rețea dintre containere ($\approx$50–100~µs) sunt mai mici
        decât într-un centru de date real, ceea ce avantajează baza de date și backend-urile la distanță.
  \item \textbf{Durată}: rulările de 15~s și repetarea unică a experimentelor HTTP introduc variații de
        câteva procente. Benchmark-ul motorului raportează mediana a 3 repetiții.
  \item \textbf{Sarcina sintetică}: modelul IRM/Zipf nu captează corelațiile temporale ale traficului real
        (de exemplu vârfuri sau sesiuni).
\end{itemize}
